LiDAR--inertial odometry (LIO) systems differ in whether they continue estimating
gravity after initialization. We compare four gravity--bias state configurations
in each of FAST-LIO2 and LIO-SAM, then separately test a gravity-direction factor.
Across \NumCoreSeq{} dataset sequences evaluated with FAST-LIO2,
fixing gravity under continuous LiDAR correction produces mean paired changes
in vertical and 3D position errors with 90\% confidence intervals within $\pm$\StrictEquivMargin{}.
Tests on \NumLioSamSeq{} sequences with LIO-SAM likewise show no consistent benefit from
online gravity. Multi-second LiDAR outages, unlike reduced range or field of
view, reveal trajectory-dependent costs of fixing gravity. A history-matched
23D-to-21D switch places the repeatable 3D error increase after LiDAR updates
resume. Under 5-s outages, a direction factor from the same IMU used for
preintegration improves accuracy on Hall05 but worsens both errors with online
gravity on TUHH. Dynamic-start tests also show fixed-bias failures at
particular starting phases. We recommend keeping gravity and accelerometer
bias online for robustness; use a direction factor only after verifying
vertical and 3D accuracy gains under the intended operating conditions.
